\par\SourceLeaf{148}\noindent
Aus den früheren Erfahrungen jedoch [S. 139 unten] hätte man unter Berücksichtigung des Umstands, dass die Gesellschaft in den ersten Jahrzehnten einen viel kleineren Umfang hatte als 1794 und dass jene Erfahrungen aus einer Zeit stammten, in der der entsprechende Beharrungszustand als noch nicht ganz erreicht angesehen werden muss, schließen können, dass man künftig eine die Durchschnittszahl 13 weit überschreitende Witwenzahl für sehr wohl möglich halten müsse. (Man darf nicht übersehen, dass bei allem, was ich hier gesagt habe, die nach \textsc{Kritter}s Gutachten wegen der Waisenpensionen noch erforderliche Vergrößerung um \(\tfrac16\) noch nicht mit einbegriffen ist.)

Die S. [140] Nr. 2 beanstandete Behauptung \textsc{Kritter}s ist zwar durch die Erfahrung hinreichend widerlegt. Es ist jedoch nicht überflüssig, zur eigentlichen, etwas versteckt liegenden Quelle des Irrtums zurückzugehen. Bei jeder Anwendung der Rechnung, sowohl auf Gegenstände der Natur als auch auf gesellschaftliche Verhältnisse, treten die Erfahrungsdaten selten in der reinen Gestalt auf, in der man sie eigentlich braucht. Fast immer sind sie mehr oder weniger mit Störungen oder Schwankungen behaftet, die in ihrem Wechsel keiner Regel folgen. Wie jeder weiß, sucht man dann den daraus entstehenden Nachteil dadurch, dass man aus vielen einzelnen Ergebnissen das Mittel nimmt, wenn schon nicht aufzuheben, so doch möglichst zu vermindern. Man rechnet damit, dass sich bei einer solchen Verwendung einer großen Zahl von Fällen die zufälligen Schwankungen größtenteils ausgleichen, und legt dem Mittelwert eine umso größere Zuverlässigkeit bei, je mehr Einzelergebnisse einbezogen werden. Das ist im Allgemeinen auch vollkommen richtig. Durch konsequente Weiterentwicklung und umsichtige Nutzung dieses Prinzips sind besonders in den Naturwissenschaften nicht selten die lohnendsten Früchte, selbst glänzende Ergebnisse, gewonnen worden. Die Sicherheit des Grundprinzips beruht jedoch auf einer wesentlichen Bedingung, die häufig genug auch von Fachgelehrten außer Acht gelassen wird: Die regellosen Störungen oder Schwankungen, die den einzelnen Beobachtungen oder Erfahrungen anhaften, müssen voneinander ganz unabhängig sein. Zu beurteilen, ob eine solche Unabhängigkeit vorliegt oder nicht, kann bisweilen sehr schwierig und ohne tiefes Eindringen in die Sachverhältnisse unmöglich sein. Bleiben darüber Zweifel bestehen, so wird auch das Gewicht, das den Endergebnissen beizulegen ist, unsicher sein.

Ginge es etwa um eine meteorologische Größe, zum Beispiel die Menge des jährlich an einem bestimmten Ort fallenden Regens, so ist diese bekanntlich in verschiedenen Jahren sehr ungleich. Den durch die allgemeinen örtlichen Verhältnisse bestimmten Normalwert wird man aber an einem Durchschnitt von zehn Jahren mit viel größerer Sicherheit erkennen, als wenn man sich nur an ein einzelnes Jahr halten wollte. Der Grund ist, dass zwischen den Abweichungen vom Normalwert in den einzelnen Jahren kein besonderer Zusammenhang besteht. Vielmehr kann, wie auch die Erfahrung bestätigt, eine große negative Abweichung ebenso leicht in einem Jahr vorkommen, das unmittelbar auf ein Jahr mit großer positiver Abweichung folgt, wie in jedem anderen.

Bei den gezählten Witwen aus aufeinanderfolgenden Jahren fehlt jedoch jene wesentliche Bedingung, gerade weil der Übergang von einer Zahl zu einer erheblich anderen nur allmählich geschehen kann. Ist zum Beispiel in der oben zur Erläuterung angeführten größeren Gesellschaft, in der der \textit{durchschnittliche} jährliche Zugang mit 69 angenommen ist und nach erreichtem Beharrungszustand der jährliche Abgang ebenso groß ist, der Bestand einmal auf 1240 gesunken oder liegt derzeit die negative Abweichung \(-60\) vor, so ist es im höchsten, an Unmöglichkeit grenzenden Grad unwahrscheinlich, dass im folgenden Jahr eine positive Abweichung vom Normalwert eintreten wird. In einer kleinen Gesellschaft wie der unsrigen sind die im folgenden Jahr gezählten Witwen sehr oft noch genau dieselben Personen wie im vorangegangenen. Selbst nach 10 Jahren wird in der Regel nur der kleinere Teil durch andere ersetzt sein. Eine Durchschnittszahl aus 10 aufeinanderfolgenden Jahren ist daher noch kein Mittel aus 10 voneinander ganz unab%
\SourceLeaf{149}%
hängigen Erfahrungen und kann über die eigentliche Normalzahl noch nicht viel sicherer Aufschluss geben als die Erfahrung eines einzelnen Jahres. Um das erhöhte Gewicht schätzen zu können, das einem Durchschnittswert beizulegen ist, kommt es wesentlich darauf an, aus wie vielen voneinander ganz unabhängigen Gruppen die Erfahrungen stammen. Da nun die durchschnittliche Dauer einer Witwenschaft etwa 20 Jahre beträgt, würde bei einer empirischen Bestimmung der Normalzahl selbst ein vierzigjähriger Durchschnitt noch keine sehr sichere Gewähr dafür bieten, dass die Schwankungen ausgeschaltet sind. Dass die Witwenzahl vierzig Jahre hintereinander ständig unter dem Normalwert bleibt, der den allgemeinen Verhältnissen der Gesellschaft entspricht, ist im Grunde ebenso wenig als eine ganz außergewöhnliche Erscheinung anzusehen, wie wenn ein Pharaospieler zweimal hintereinander gewinnt. Die Lehre, die man daraus ziehen muss, ist, dass umgekehrt in dieser Hinsicht auch 40 ununterbrochen magere Jahre ebenso leicht möglich sind wie 40 ununterbrochen fette.

